本组采用基础技术方向与扩展模块相结合的分工方式。四名成员分别承担前端、后端、数据库和测试的主要职责，同时负责看板、红包积分、AI 与高德地图等扩展内容。表\ref{tab:team-process}依据小组确认的实际分工整理，不以模块数量推定工作量比例。
\begin{longtable}{L{.15\linewidth}L{.16\linewidth}L{.16\linewidth}L{.42\linewidth}}
\caption{小组实际分工}\label{tab:team-process}\\\toprule 成员&基础方向&扩展模块&协作接口与联调关注点\\\midrule\endfirsthead\toprule 成员&基础方向&扩展模块&协作接口与联调关注点\\\midrule\endhead
范文丽&前端&看板&对齐后端响应与页面状态，明确看板统计口径和数据展示。\\
李承文&后端&红包积分&对齐订单、优惠券与积分规则，关注计价、支付和取消回退。\\
赵紫怡&数据库&AI&协调业务数据与查询，关注真实菜单、订单授权和 AI 调用边界。\\
仲禹潼&测试&高德地图&组织功能与回归验证，关注定位、导航及地图失败后的替代路径。\\\bottomrule
\end{longtable}
这一安排结合了成员的技术主责与具体功能任务，但也意味着一个需求往往需要多人共同完成。以 AI 点餐为例，候选商品依赖数据库与 AI 模块，展示和确认依赖前端，加入购物车仍需后端执行规则，最后由测试检查完整流程。高德地图同样需要页面入口、后端数据与异常处理共同配合。

因此，测试由专人主责不意味着其他成员可以不验证自己的代码，前端主责也不意味着只需要了解页面。结合答辩意见，我们需要在保留技术主责的同时，加强对完整业务链路的理解，让每项需求都有明确的联调与验收责任。

在协作组织上，小组使用腾讯会议进行阶段性正式讨论，使用微信群处理即时问题，使用飞书共享文档维护当前任务、原始记录、总结材料和教师要求，并通过 GitHub 分支与 Pull Request 完成代码协作。普通项目文档原则上保存在仓库的 \code{docs} 目录中；外部共享文档的访问地址记录在仓库根目录的 \code{README.md} 中，以便成员在不同开发阶段查阅。成员分工表描述的是主要责任范围，而不是文件或代码目录的排他所有权；涉及跨层业务时，由相关成员共同确认接口、数据、页面状态和验收条件。
